\section{Additional Considerations, Abbreviations, and Amendment History}

\subsection{Additional Considerations}

Although this protocol is written to the United States combined IND/IDE pathway, the
LLM-directed robotic system is a Software as a Medical Device (SaMD) whose
deployment may be sought across jurisdictions, and the design anticipates the
parallel regulatory tracks summarized in Table~\ref{tab:jurisdictions}. The European
Union Medical Device Regulation (EU MDR) and its companion artificial-intelligence
requirements, China's National Medical Products Administration (NMPA), Japan's
Pharmaceuticals and Medical Devices Agency (PMDA), the United Kingdom's Medicines
and Healthcare products Regulatory Agency (MHRA), Health Canada, and Australia's
Therapeutic Goods Administration (TGA) each maintain a SaMD and robotically assisted
surgery framework with which the device evidence is designed to remain compatible.
The common technical spine, Phase 0 simulation validation, USL readiness rating, the
federated hash-chained audit trail, and the verification-before-generation gate
surface, is constructed to map onto each track without redesign, and the
international consensus standards cited in $\S$12 (IEC 80601-2-77, ASME V and V 40,
NASA-STD-7009A, IEEE 7009-2024) provide a shared basis across jurisdictions
\cite{iec8060127,asmevv40,nasastd7009,ieee7009}.

A multicenter Physical AI trial also depends on learning across sites without
centralizing raw patient data. The eight academic HPB centers operate as a federated
network: model and validation evidence is exchanged across sites while raw telemetry
and identifiable records remain local, the audit trail is hash-chained across the
fleet so that every site's record is tamper-evident, and fleet harmonization ensures
that every device instance behaves identically. This federated design lets the
network accumulate validation evidence and reproduce simulated procedures at scale
while preserving the deny-by-default privacy model of $\S$10.6 and the upgraded
Phase 0 readiness gate at each site.

The trial is additionally aligned with the legislative and financial framework of
the Verification Before Generation in Physical AI Oncology Trials Act of 2026 (H.R.
9510) \cite{kawchak2026hr9510}, which ties the funding and oversight of Physical AI
oncology investigations to a verification-and-validation-and-uncertainty-quantification
(VVUQ) evidentiary standard. That framework supplies the financial-data standard
under which the Patient-Aligned Co-Investment Facility operates: co-investment
capital raises operational capacity and trial success likelihood while the capital
firewall and the VVUQ standard together ensure that funding is verifiable and that
no funder can reach the science of the comparison. The ten-gate assurance suite
operationalizes the technical half of that standard: every proposed robot-patient
command is verified before it is generated, against fourteen external safety and
robotics standards plus two clinical baselines, with epistemic and aleatory
uncertainty bounds, and the gate surface returns one of accept, block, or escalate.
Figure 22 renders the ten-gate pipeline.

\begin{table}[htbp]
\centering
\footnotesize
\caption{Cross-jurisdictional Software as a Medical Device (SaMD) and robotically
assisted surgery tracks. The common technical spine (Phase 0 validation, USL
readiness, the federated hash-chained audit trail, and the
verification-before-generation gate surface) maps onto each track without redesign.}
\label{tab:jurisdictions}
\begin{tabularx}{\textwidth}{L{2.5cm} L{3.2cm} Y}
\toprule
\textbf{Jurisdiction} & \textbf{Regulator / framework} & \textbf{Alignment basis} \\
\midrule
European Union & EU MDR + AI reqs & SaMD risk classification; common simulation-validation and federated audit evidence. \\
China & NMPA & SaMD, robotically assisted surgery review; standards-based VVUQ. \\
Japan & PMDA & SaMD and device review; consensus-standard alignment. \\
United Kingdom & MHRA & SaMD framework; post-market and audit-trail compatibility. \\
Canada & Health Canada & SaMD licensing; quality-system and validation evidence. \\
Australia & TGA & SaMD inclusion; standards-based safety and performance evidence. \\
\bottomrule
\end{tabularx}
\end{table}

\vspace{-0.2cm}

\begin{mermaidfig}
\node[mmin]  (gen)  at (-0.2,0)      {Proposed robot-patient\\interaction code /\\command};
\node[mmstep](gates)at (4.6,0)    {10-gate assurance suite\\14 external standards\\+ 2 clinical baselines};
\node[mmstep](uq)   at (9.6,0)    {Uncertainty\\quantification:\\epistemic + aleatory};
\node[mmdec] (v)    at (9.6,-2.6) {Gate surface\\verdict};
\node[mmgoal](acc)  at (4.6,-2.6) {ACCEPT: verified\\before generation};
\node[mmdark](blk)  at (-0.2,-2.6)   {BLOCK: command\\refused};
\node[mmdark](esc)  at (4.6,-4.5) {ESCALATE: hand\\back to human};
\draw[mmedge] (gen) -- (gates);
\draw[mmedge] (gates) -- (uq);
\draw[mmedge] (uq) -- (v);
\draw[mmedge] (v) -- node[above,font=\scriptsize, xshift=0.1cm]{accept} (acc);
\draw[mmedge] (acc) -- node[above,font=\scriptsize, xshift=0.05cm]{block} (blk);
\draw[mmedge] (v.south west) -- node[right,font=\scriptsize,pos=0.6]{escalate} (esc.east);
\end{mermaidfig}
\vspace{-0.7cm}
\figcaption{Figure 22. Verification before generation. Every proposed robot-patient
command passes a ten-gate assurance suite aligned to fourteen external safety and
robotics standards plus two clinical baselines, with epistemic and aleatory
uncertainty bounds, before the gate surface returns ACCEPT, BLOCK, or ESCALATE (hand
back to human).}

\subsection{Protocol Amendment History}

Table~\ref{tab:amend} records the version history of this protocol and its Phase 1
predicate; subsequent amendments will be appended with the version, date, sections
affected, and a summary of the change.

\begin{table}[htbp]
\centering
\footnotesize
\caption{Protocol amendment history.}
\label{tab:amend}
\begin{tabularx}{\textwidth}{L{2.0cm} L{2.6cm} L{3.2cm} Y}
\toprule
\textbf{Version} & \textbf{Date} & \textbf{Sections affected} & \textbf{Summary of change} \\
\midrule
v1.0.0 & June 20, 2026 & All (predicate protocol) & Phase 1 predicate protocol \cite{kawchak2026phase1}: the first-in-human, single-arm, combined IND/IDE design that established the daraxonrasib recommended Phase 2 dose through a perioperative 3+3, the on-premises LLM-directed eight-arm robotic Whipple device feasibility, the Physical AI Subpart J overlay, and the safety and oversight framework on which this Phase 2 builds. \\
v1.1.0 & June 23, 2026 & All (this protocol) & This Phase 2, multicenter, randomized, controlled protocol. Establishes the randomized 1:1 parallel-group design with BICR, daraxonrasib at the recommended Phase 2 dose plus standard-of-care control, the upgraded Phase 0 readiness gate (USL $\geq 8.0$), multicenter single-IRB governance, the co-investment model and the capital firewall, and the confirmatory group-sequential statistical framework. \\
\bottomrule
\end{tabularx}
\end{table}

\vspace{-0.2cm}

\subsection{Abbreviations}

Table~\ref{tab:abbrev} resolves every abbreviation used across this protocol.

\begin{table}[htbp]
\centering
\footnotesize
\caption{Abbreviations used in this protocol.}
\label{tab:abbrev}
\begin{tabularx}{\textwidth}{L{1.3cm} L{6.8cm} L{1.3cm} L{6.2cm}}
\toprule
\textbf{Abbr.} & \textbf{Definition} & \textbf{Abbr.} & \textbf{Definition} \\
\midrule
AE     & Adverse event                                  & MTD   & Maximum tolerated dose \\
BICR   & Blinded independent central review              & NMPA  & National Medical Products Administration \\
CONSORT& Consolidated Standards of Reporting Trials      & OS    & Overall survival \\
CRF    & Case report form                                & PACIF & Patient-Aligned Co-Investment Facility \\
CRO    & Clinical research organization                  & PDAC  & Pancreatic ductal adenocarcinoma \\
ctDNA  & Circulating tumor DNA                           & PFS   & Progression-free survival \\
CTCAE  & Common Terminology Criteria for Adverse Events  & PJ    & Pancreaticojejunostomy \\
DLT    & Dose-limiting toxicity                          & PMDA  & Pharmaceuticals and Medical Devices Agency \\
DSMB   & Data and Safety Monitoring Board                & PP    & Per-protocol \\
ECOG   & Eastern Cooperative Oncology Group              & PV    & Portal vein \\
EORTC  & European Organisation for Research and Treatment of Cancer & QoL & Quality of life \\
FDA    & U.S. Food and Drug Administration               & R0    & Margin-negative resection \\
GCP    & Good Clinical Practice                          & RCT   & Randomized controlled trial \\
GJ     & Gastrojejunostomy                               & RP2D  & Recommended Phase 2 dose \\
HJ     & Hepaticojejunostomy                             & SAE   & Serious adverse event \\
HPB    & Hepatopancreatobiliary                          & SaMD  & Software as a Medical Device \\
HR     & Hazard ratio                                    & sIRB  & Single Institutional Review Board \\
ICH    & International Council for Harmonisation         & SMA   & Superior mesenteric artery \\
IDE    & Investigational Device Exemption                & SMV   & Superior mesenteric vein \\
IND    & Investigational New Drug                        & TGA   & Therapeutic Goods Administration \\
ISGPS  & International Study Group of Pancreatic Surgery & USL   & Unified Safety Level \\
ISM    & Independent Safety Monitor                      & VVUQ  & Verification, validation, and uncertainty quantification \\
ITT    & Intention-to-treat                              & MDR   & Medical Device Regulation \\
KRAS   & Kirsten rat sarcoma viral oncogene homolog      & MHRA  & Medicines and Healthcare products Regulatory Agency \\
LLM    & Large language model                            & MPR   & Major pathologic response \\
mITT   & Modified intention-to-treat                     &       & \\
\bottomrule
\end{tabularx}
\end{table}
